\slajdid{02-s041}
\begin{frame}[label=02-s041]
\gusttekst
Iz funkcionalne tabele, ili iz specifikacije kombinacione mreže, se u Karnoovu tabelu unose vrednosti funkcije za stanja ulaznih promenljivih u odgovarajuće polje.

Algoritam minimizacije
\begin{enumerate}
\item Da bi dobili funkciju u obliku zbira proizvoda potrebno je sve logičke jedinice koje se pojavljuju u Karnoovoj tabeli prekriti sa \alert{što manjim brojem površina što višeg reda}.
\item Da bi dobili funkciju u obliku proizvoda zbirova potrebno je sve logičke nule koje se pojavljuju u Karnoovoj tabeli prekriti sa \alert{što manjim brojem površina što višeg reda}.
\end{enumerate}
Ako je ispunjen uslov da je broj površina minimalan sa najvećim mogućim površinama, algoritam garantuje da će i dobijene funkcije biti minimalne i da minimalnije od toga nije moguće realizovati funkcije korišćenjem I i ILI logičkih kola, podrazumevajući da raspolažemo i sa pravim i komplementnim vrednostima signala na ulazu u kombinacionu mrežu.
\end{frame}
